\section{Why Biomolecular AI Is a Particularly Strong Testbed}

The original paper already pointed toward structural biology as a natural application area \citep{corcoran2018spatial}. That instinct looks stronger now than it did in 2018. Biomolecular learning combines nearly all of the conditions under which serialization becomes attractive: sparse but information-dense 3D structure, strong local motifs, important long-range dependencies, and a need to carry multiple semantic channels per structural element.

Recent work such as AtoMAE confirms that voxel-based biomolecular representation learning remains scientifically active when paired with modern architectures \citep{park2025atomae}. This does not mean that one representation should dominate all others. Rather, it suggests that biomolecular AI is still searching for the right interface between structure and model family. Sequence models are natural for proteins because biological data already come with sequences. Structural models are natural because function depends on three-dimensional arrangement. Serialization matters because it creates one of the few practical bridges between those two worlds.

From this viewpoint, a useful biomolecular serialization strategy should satisfy four design criteria.

\paragraph{Locality preservation.}
The ordering should preserve enough spatial coherence that local motifs remain learnable without massive receptive fields. This is the most direct continuity with the 2018 Hilbert-mapping argument.

\paragraph{Token efficiency.}
The representation must control token count and neighborhood access so that structure can be processed at realistic biological scales. Serialization becomes especially valuable when it turns expensive geometric search into cheap ordered lookup or sequence processing.

\paragraph{Channel semantics.}
Each token should be able to carry chemically meaningful information. A serialization that preserves geometry but discards atom and residue semantics will underperform on scientific tasks whose labels depend on more than shape.

\paragraph{Reversibility and interpretability.}
Scientific users often need to map learned evidence back to physical structure. Reversible or at least traceable orderings are therefore more valuable than opaque latent compressions. The source paper's emphasis on reversible mappings was ahead of its time in this respect.

These criteria suggest a practical shift in how future papers should be written. A biomolecular serialization paper should not be evaluated solely as a backbone paper. It should also be evaluated as an interface paper: a paper about what information is retained, what gets reordered, what becomes cheaply accessible, and what explanations remain possible after tokenization.
